% Sub-section: Hazard
\begin{frame}{3.3 Quản Lý Xung Đột (\texttt{pipeline\_hazard.v})}
    Chạy gối đầu lệnh tạo ra các xung đột phần cứng nguy hiểm cần mạch Logic xử lý:
    \vspace{0.3cm}
    \begin{columns}
        \begin{column}{0.5\textwidth}
            \textbf{1. Data Hazard (Xung đột dữ liệu)}
            \begin{itemize}
                \item Xảy ra khi lệnh sau cần đọc thanh ghi mà lệnh trước chưa kịp ghi xong.
                \item \textbf{Giải pháp:} Thiết kế mạch \textbf{Forwarding Unit} để bơm tắt dữ liệu từ ngõ ra ALU tầng sau về ngay ngõ vào ALU tầng trước mà không cần đợi.
                \item Với lệnh Load, thực hiện \textbf{Stall} (hoãn) pipeline 1 chu kỳ.
            \end{itemize}
        \end{column}
        
        \begin{column}{0.5\textwidth}
            \textbf{2. Control Hazard (Xung đột điều khiển)}
            \begin{itemize}
                \item Xảy ra khi lệnh rẽ nhánh (\texttt{BEQ}, \texttt{BNE}) làm thay đổi dòng chảy của \texttt{PC}.
                \item \textbf{Giải pháp:} Thực hiện \textbf{Flush} (xóa bỏ dữ liệu lỗi bằng cách nạp mã lệnh rỗng NOP) tại các tầng nạp sai lệnh trước đó.
            \end{itemize}
        \end{column}
    \end{columns}
\end{frame}